\documentclass[10pt,a4paper]{article}
\usepackage[margin=2cm]{geometry}
\usepackage[T1]{fontenc}
\usepackage{lmodern}
\usepackage{amsmath,amssymb,booktabs,multirow,xcolor,colortbl}
\usepackage{pgfplots}
\pgfplotsset{compat=1.18}
\usepackage{float}
\usepackage[hidelinks]{hyperref}
\usepackage{titlesec,enumitem}
\titleformat{\section}{\large\bfseries}{\thesection}{0.6em}{}
\setlist{itemsep=1pt,topsep=3pt}

\definecolor{cgrad}{HTML}{C0612B}
\definecolor{crand}{HTML}{8A939C}
\pgfplotsset{
  accplot/.style={width=0.49\textwidth,height=5.6cm,grid=major,grid style={gray!20},
    tick label style={font=\scriptsize},label style={font=\footnotesize},title style={font=\small\bfseries},
    legend style={font=\scriptsize,draw=none,fill=white,fill opacity=0.85,text opacity=1},
    legend cell align=left,mark size=1.3pt},
  pslc/.style={black,densely dotted,thick,mark=square*,mark options={solid}},
  pmlc/.style={black!55,dashdotted,thick,mark=triangle*,mark options={solid}},
  pgrad/.style={cgrad,very thick,mark=*},
  prand/.style={crand,thick,mark=o}}

\begin{document}

\begin{center}
{\LARGE\bfseries Baseline Results}\\[4pt]
{\large Reliable Hybrid SLC/MLC ReRAM Crossbars for Transformers}\\[4pt]
Shivansh Mittal, Sarvesh Joshi, Thumati Naveen Chandra Chowdary, Aryan Chakravorty\\
CS5290 Memory-Centric Computer Architecture, Indian Institute of Technology Guwahati
\end{center}

\section{What the baseline is}
The baseline is \textbf{HyFlexPIM} (Song \emph{et al.}, ISCA 2025), reproduced on our cell-level crossbar
simulator. Every encoder linear layer of BERT-Base (72 matrices) and DistilBERT (36 matrices) is factorised
with SVD, keeping $k=mn/(m+n)$ components so that storage is unchanged, and fine-tuned on SST-2 for 8
epochs. Components are ranked by the HyFlexPIM gradient score, $\sum_b|\partial L_b/\partial\sigma_r|$ over
batches of 32 training sentences, and the top 20\% of each layer go to SLC cells, the rest to MLC cells.
Three reference chips are evaluated alongside it:
\begin{itemize}
\item \textbf{All SLC}: every component in reliable 1-bit cells (14 cells per weight), the most robust and most
expensive chip;
\item \textbf{All MLC}: every component in dense 2-bit cells (8 cells per weight), the cheapest chip;
\item \textbf{Random 20\%}: a random 20\% of components in SLC, a control for whether the choice matters.
\end{itemize}
The hybrid chips use 9.2 cells per weight, 34\% fewer than all-SLC. Accuracy is measured on all 872 SST-2
validation sentences and averaged over several simulated chips (random draws of all cell errors). Every
placement is tested on the same chips. Errors are modelled per cell: 8-bit weights with one scale per output
column, differential cell pairs, 7 SLC or 4 MLC bit slices per sign, and programming noise
$\tilde G=G+\sigma G_{\max}\varepsilon$ on every cell.

\section{Accuracy before hardware errors}
\begin{table}[H]
\centering\small
\caption{SST-2 accuracy (\%) without hardware errors.}
\begin{tabular}{@{}lcccc@{}}
\toprule
Model & FP32 & INT8 & SVD, truncated & SVD, fine-tuned \\
\midrule
BERT-Base & 92.43 & 91.86 & 89.33 & 90.14 \\
DistilBERT & 91.06 & 90.37 & 89.68 & 89.11 \\
\bottomrule
\end{tabular}
\end{table}
\noindent The SVD factorisation costs about two points after fine-tuning; this cost is common to every
placement. On an error-free simulated chip every placement stays within about one point of the digital
SVD model, confirming the simulator.

\section{Programming noise}
\begin{figure}[H]
\centering
\begin{tikzpicture}
\begin{axis}[accplot,title={BERT-Base},xlabel={Programming noise $\sigma$ (\% of $G_{\max}$)},ylabel={Accuracy (\%)},
  xtick={0,1,2,3,5,7.5,10},ymin=45,ymax=93,legend pos=south west]
\addplot[pslc] table[x=x,y=all_slc]{data/noise_bert.dat}; \addlegendentry{All SLC}
\addplot[pgrad] table[x=x,y=gradient]{data/noise_bert.dat}; \addlegendentry{HyFlexPIM (gradient) 20\%}
\addplot[prand] table[x=x,y=random]{data/noise_bert.dat}; \addlegendentry{Random 20\%}
\addplot[pmlc] table[x=x,y=all_mlc]{data/noise_bert.dat}; \addlegendentry{All MLC}
\end{axis}
\end{tikzpicture}\hfill
\begin{tikzpicture}
\begin{axis}[accplot,title={DistilBERT},xlabel={Programming noise $\sigma$ (\% of $G_{\max}$)},
  xtick={0,1,2,3,5,7.5,10},ymin=45,ymax=93]
\addplot[pslc] table[x=x,y=all_slc]{data/noise_distil.dat};
\addplot[pgrad] table[x=x,y=gradient]{data/noise_distil.dat};
\addplot[prand] table[x=x,y=random]{data/noise_distil.dat};
\addplot[pmlc] table[x=x,y=all_mlc]{data/noise_distil.dat};
\end{axis}
\end{tikzpicture}
\caption{Accuracy against programming noise, mean of 5 chips.}
\end{figure}

\begin{table}[H]
\centering\small
\caption{Accuracy (\%) against programming noise $\sigma$, mean of 5 chips.}
\begin{tabular}{@{}llccccccc@{}}
\toprule
Model & Placement & 0 & 1\% & 2\% & 3\% & 5\% & 7.5\% & 10\% \\
\midrule
\multirow{4}{*}{BERT-Base} & All SLC & 89.9 & 89.6 & 89.3 & 89.1 & 87.8 & 85.9 & 79.7 \\
 & HyFlexPIM (gradient) 20\% & 89.8 & 90.4 & 90.2 & 90.0 & 88.0 & 79.3 & 67.7 \\
 & Random 20\% & 90.0 & 90.1 & 89.9 & 87.7 & 78.8 & 58.8 & 49.1 \\
 & All MLC & 90.0 & 89.3 & 87.3 & 83.9 & 66.3 & 52.2 & 51.3 \\
\midrule
\multirow{4}{*}{DistilBERT} & All SLC & 89.0 & 88.7 & 88.8 & 88.6 & 88.5 & 87.5 & 85.9 \\
 & HyFlexPIM (gradient) 20\% & 89.8 & 88.5 & 88.9 & 88.7 & 87.7 & 83.4 & 69.8 \\
 & Random 20\% & 88.6 & 88.7 & 88.7 & 87.8 & 82.8 & 61.8 & 52.4 \\
 & All MLC & 89.1 & 88.4 & 87.9 & 86.1 & 74.5 & 60.4 & 56.0 \\
\bottomrule
\end{tabular}
\end{table}

\begin{itemize}
\item \textbf{All-MLC fails at moderate noise}: BERT-Base falls to 66.3\% at $\sigma=5\%$ and to near chance
(52.2\%) at 7.5\%. A hybrid layout is necessary.
\item \textbf{The HyFlexPIM hybrid works at low noise but degrades at high noise}: with 20\% SLC it matches
all-SLC up to $\sigma=5\%$, then trails it by 6.6 points at 7.5\% and 12 points at 10\% on BERT-Base.
\item \textbf{Random selection recovers little}: which components are protected matters, not only how many.
\end{itemize}

\section{SLC share}
\begin{figure}[H]
\centering
\begin{tikzpicture}
\begin{axis}[accplot,title={BERT-Base, $\sigma=7.5\%$},xlabel={Share of components in SLC (\%)},ylabel={Accuracy (\%)},
  xtick={0,10,20,30,50,100},ymin=50,ymax=90,legend pos=south east]
\addplot[pgrad] table[x=x,y=gradient]{data/share_bert_075.dat}; \addlegendentry{HyFlexPIM (gradient)}
\addplot[prand] table[x=x,y=random]{data/share_bert_075.dat}; \addlegendentry{Random}
\end{axis}
\end{tikzpicture}\hfill
\begin{tikzpicture}
\begin{axis}[accplot,title={DistilBERT, $\sigma=7.5\%$},xlabel={Share of components in SLC (\%)},
  xtick={0,10,20,30,50,100},ymin=50,ymax=90]
\addplot[pgrad] table[x=x,y=gradient]{data/share_distil_075.dat};
\addplot[prand] table[x=x,y=random]{data/share_distil_075.dat};
\end{axis}
\end{tikzpicture}
\caption{Accuracy against the share of components in SLC at $\sigma=7.5\%$, mean of 20 chips; 0\% is all-MLC
and 100\% all-SLC.}
\end{figure}

\begin{table}[H]
\centering\small
\caption{Accuracy (\%) of the HyFlexPIM gradient placement by SLC share, mean of 20 chips.}
\begin{tabular}{@{}llccccccc@{}}
\toprule
Model & $\sigma$ & 0\% (all MLC) & 5\% & 10\% & 20\% & 30\% & 50\% & 100\% (all SLC) \\
\midrule
\multirow{3}{*}{BERT-Base} & 3\% & 82.1 & 88.7 & 90.0 & 90.0 & 89.9 & 89.6 & 89.2 \\
 & 5\% & 66.8 & 84.2 & 88.2 & 88.4 & 88.4 & 88.1 & 88.0 \\
 & 7.5\% & 53.0 & 67.5 & 79.0 & 81.2 & 82.4 & 83.2 & 85.7 \\
\midrule
\multirow{3}{*}{DistilBERT} & 3\% & 86.0 & 88.4 & 88.5 & 88.6 & 88.8 & 88.7 & 88.7 \\
 & 5\% & 74.8 & 87.3 & 87.3 & 87.7 & 88.1 & 87.8 & 88.2 \\
 & 7.5\% & 56.9 & 78.0 & 80.7 & 83.1 & 84.8 & 86.0 & 87.4 \\
\bottomrule
\end{tabular}
\end{table}

\noindent At $\sigma=7.5\%$ the gradient placement does not reach all-SLC even with 50\% of the components in
SLC on BERT-Base (83.2\% against 85.7\%), so most of the cell savings of the hybrid are lost at high noise.

\section{Stuck cells}
\begin{figure}[H]
\centering
\begin{tikzpicture}
\begin{axis}[accplot,title={BERT-Base},xlabel={Stuck cells $p$ (\% of cells)},ylabel={Accuracy (\%)},
  symbolic x coords={0,0.1,0.2,0.5,1,2,5},xtick=data,ymin=45,ymax=93,legend pos=north east]
\addplot[pslc] table[x=x,y=all_slc]{data/fault_bert.dat}; \addlegendentry{All SLC}
\addplot[pgrad] table[x=x,y=gradient]{data/fault_bert.dat}; \addlegendentry{HyFlexPIM 20\%}
\addplot[prand] table[x=x,y=random]{data/fault_bert.dat}; \addlegendentry{Random 20\%}
\addplot[pmlc] table[x=x,y=all_mlc]{data/fault_bert.dat}; \addlegendentry{All MLC}
\end{axis}
\end{tikzpicture}\hfill
\begin{tikzpicture}
\begin{axis}[accplot,title={DistilBERT},xlabel={Stuck cells $p$ (\% of cells)},
  symbolic x coords={0,0.1,0.2,0.5,1,2,5},xtick=data,ymin=45,ymax=93]
\addplot[pslc] table[x=x,y=all_slc]{data/fault_distil.dat};
\addplot[pgrad] table[x=x,y=gradient]{data/fault_distil.dat};
\addplot[prand] table[x=x,y=random]{data/fault_distil.dat};
\addplot[pmlc] table[x=x,y=all_mlc]{data/fault_distil.dat};
\end{axis}
\end{tikzpicture}
\caption{Accuracy against the share of stuck cells (half at $G_{\max}$, half at $G_{\min}$), no programming
noise, mean of 20 chips. The horizontal axis is not to scale.}
\end{figure}

\begin{itemize}
\item \textbf{Stuck cells break every baseline}: all-MLC fails at 0.1\% stuck cells (58.6\% on BERT-Base),
the HyFlexPIM hybrid drops to 71.0\% at 0.5\% and 52.8\% at 1\%, and even all-SLC falls to 62.7\% at 1\%.
\item A stuck cell in a high-order slice shifts its weight by a large fixed amount, so faults are far more
harmful than noise, and accuracy varies widely between chips depending on where the faults fall.
\end{itemize}

\section{Other non-idealities}
\begin{table}[H]
\centering\small
\caption{Accuracy (\%) under each other non-ideality, one at a time. The hybrid keeps 20\% of the components
in SLC by the HyFlexPIM gradient score.}
\setlength{\tabcolsep}{5pt}
\begin{tabular}{@{}lccc@{\hskip 14pt}ccc@{}}
\toprule
& \multicolumn{3}{c}{BERT-Base} & \multicolumn{3}{c}{DistilBERT} \\
\cmidrule(r){2-4}\cmidrule(l){5-7}
Condition & All SLC & All MLC & HyFlexPIM & All SLC & All MLC & HyFlexPIM \\
\midrule
No errors & 89.7 & 90.4 & 90.3 & 88.5 & 89.0 & 88.6 \\
Read noise $\sigma_r=2\%$ & 89.0 & 77.3 & 89.8 & 88.3 & 83.8 & 88.1 \\
Read noise $\sigma_r=5\%$ & 80.0 & 53.0 & 76.6 & 84.5 & 57.7 & 80.5 \\
$\sigma=5\%$, middle MLC levels $\times2$ & 88.2 & 62.5 & 87.8 & 88.4 & 70.3 & 87.1 \\
1\% stuck, scattered & 61.3 & 50.7 & 52.6 & 74.2 & 51.2 & 55.0 \\
1\% stuck, $3\times3$ clusters & 66.7 & 50.8 & 51.5 & 80.3 & 51.2 & 56.4 \\
1\% stuck, bit lines & 57.0 & 50.4 & 51.0 & 71.2 & 50.8 & 52.2 \\
ADC 7 bits & 90.0 & 89.2 & 90.1 & 88.4 & 87.8 & 88.5 \\
ADC 6 bits & 89.3 & 83.0 & 89.6 & 88.0 & 81.7 & 88.4 \\
ADC 5 bits & 51.1 & 51.1 & 53.0 & 71.0 & 56.5 & 77.3 \\
IR drop $r=3\times10^{-4}$ & 88.9 & 89.8 & 90.5 & 89.4 & 89.1 & 89.6 \\
IR drop $r=3\times10^{-3}$ & 58.5 & 79.1 & 76.4 & 70.1 & 77.9 & 74.7 \\
\bottomrule
\end{tabular}
\end{table}

\begin{itemize}
\item Read noise and level-dependent noise behave like programming noise: all-MLC fails, and the HyFlexPIM
hybrid loses 3.4--4 points against all-SLC at 5\% read noise.
\item A fixed 6-bit ADC is enough for all-SLC and the hybrid; all-MLC needs 7 bits; 5 bits fails BERT-Base.
\item IR drop has no effect at moderate wire resistance and hurts SLC most at high resistance, because a
stored SLC one sits at the highest conductance.
\item Every fault pattern leaves the HyFlexPIM hybrid close to chance at 1\% stuck cells.
\end{itemize}

\section{Second dataset: MNLI}
The same baselines on MNLI (three classes, chance 33.3\%), on 2{,}000 matched validation pairs. On MNLI the
simulator loses accuracy even on an error-free chip (up to 11 points for all-SLC on BERT-Base) because the
vector between a tier's two crossbars is re-quantised with a single 8-bit scale; only placements with the same
tier sizes should therefore be compared directly.

\begin{table}[H]
\centering\small
\caption{MNLI accuracy (\%) of the baselines.}
\setlength{\tabcolsep}{4pt}
\begin{tabular}{@{}lcccccc@{}}
\toprule
& \multicolumn{2}{c}{$\sigma=5\%$} & \multicolumn{4}{c}{Stuck cells (no noise)} \\
\cmidrule(r){2-3}\cmidrule(l){4-7}
Placement & BERT & DistilBERT & BERT 0.1\% & BERT 0.5\% & DistilBERT 0.1\% & DistilBERT 0.5\% \\
\midrule
Digital SVD model & 83.2 & 80.9 & 83.2 & 83.2 & 80.9 & 80.9 \\
All SLC & 61.6 & 72.0 & 67.2 & 44.4 & 74.4 & 60.4 \\
HyFlexPIM (gradient) 20\% & 63.3 & 71.3 & 73.0 & 36.9 & 76.0 & 44.5 \\
Random 20\% & 35.9 & 44.1 & 42.8 & 33.3 & 51.6 & 32.9 \\
All MLC & 33.5 & 36.8 & 37.2 & 32.8 & 44.5 & 33.6 \\
\bottomrule
\end{tabular}
\end{table}

\section{Summary of the baseline}
\begin{itemize}
\item A hybrid SLC/MLC layout is necessary: all-MLC fails at moderate programming noise and at 0.1\% stuck
cells.
\item The HyFlexPIM gradient placement works at low noise but falls behind all-SLC at $\sigma\geq7.5\%$,
by 4--16 points at 20\% SLC and by up to 18 points at 5\% SLC on BERT-Base.
\item No baseline handles stuck cells: at 1\% stuck cells the HyFlexPIM hybrid is near chance and all-SLC
loses 14--27 points.
\end{itemize}
These gaps are what our contributions address; see the separate \emph{Contribution Results} document.

\end{document}
